
\documentclass[11pt,a4paper]{article}

\usepackage[a4paper,margin=2.55cm]{geometry}
\usepackage[utf8]{inputenc}
\usepackage[T1]{fontenc}
\usepackage{amsmath,amssymb,amsthm,mathtools}
\usepackage{booktabs}
\usepackage{array}
\usepackage{microtype}
\usepackage{hyperref}
\usepackage{xcolor}
\usepackage{enumitem}
\usepackage{bm}

\hypersetup{
    colorlinks=true,
    linkcolor=blue,
    citecolor=blue,
    urlcolor=blue
}

\newtheorem{theorem}{Theorem}[section]
\newtheorem{proposition}[theorem]{Proposition}
\newtheorem{lemma}[theorem]{Lemma}
\newtheorem{corollary}[theorem]{Corollary}
\newtheorem{definition}[theorem]{Definition}
\newtheorem{remark}[theorem]{Remark}
\newtheorem{problem}[theorem]{Open Problem}

\newcommand{\Hh}{\mathrm H}
\newcommand{\I}{\mathrm I}
\newcommand{\Law}{\operatorname{Law}}
\newcommand{\MEC}{\operatorname{MEC}}
\newcommand{\spec}{\operatorname{spec}}
\newcommand{\supp}{\operatorname{supp}}
\newcommand{\TC}{\operatorname{TC}}
\newcommand{\Var}{\operatorname{Var}}
\newcommand{\Cov}{\operatorname{Cov}}
\newcommand{\Csyn}{C^{\mathrm{syn}}}
\newcommand{\Ccausal}{C^{\mathrm{causal}}}
\newcommand{\CinfSyn}{C_{\infty}^{\mathrm{syn}}}
\newcommand{\CinfCausal}{C_{\infty}^{\mathrm{causal}}}
\newcommand{\cA}{\mathcal A}
\newcommand{\cY}{\mathcal Y}
\newcommand{\cM}{\mathcal M}
\newcommand{\cH}{\mathcal H}
\newcommand{\cO}{\mathcal O}
\newcommand{\cG}{\mathcal G}
\newcommand{\E}{\mathbb E}
\newcommand{\R}{\mathbb R}
\newcommand{\C}{\mathbb C}
\newcommand{\1}{\mathbf 1}
\newcommand{\dd}{\,\mathrm d}
\newcommand{\ip}[2]{\langle #1,#2\rangle}

\title{
\textbf{Who Pays the Bill?}\\[1mm]
\large
Counterfactual Identity, the Causal Shannon Floor,\\
Spectral Ledgers, and Orbit Resolution
}

\author{Jan Mikulík}
\date{20 August 2026\\[1mm]\textbf{Version 5}}

\begin{document}
\maketitle

% ============================================================
\begin{abstract}
% ============================================================

We study the minimum Shannon entropy required by exact latent
realizations of a finite memoryless controlled channel once the
counterfactual realization contract is specified.

Four contracts are separated.
Chosen-path sampling requires only the output actually requested.
A one-step counterfactual table jointly defines the outputs of all
actions at one time.
The synchronous stable-event contract C3-S requires a single
time/action coordinate to retain the same latent identity across
mutually exclusive histories.
The general causal contract C3-C allows the corresponding latent event
to depend on the realized history.

For C3-S we retain the synchronous block theory developed in earlier
versions.
The block cost is subadditive, hidden cross-counterfactual synergy can
strictly lower the naive repeated one-shot cost, and for two actions
\(P,Q\) one has the information-spectrum converse
\[
C_n^{\mathrm{syn}}(P,Q)
\ge
\frac n2
\left[
H(P)+H(Q)+W_1(\mu_P,\mu_Q)
\right],
\]
where \(\mu_P,\mu_Q\) are the laws of one-symbol surprisal.
For the equal-entropy three-outcome crash test,
\[
1.728465468745531
\le
C_\infty^{\mathrm{syn}}
\le
1.939147788079282.
\]

For C3-C we prove a different asymptotic theorem.
For every finite two-action memoryless controlled channel,
\[
\boxed{
C_\infty^{\mathrm{causal}}(P,Q)
=
\max\{H(P),H(Q)\}.
}
\]
The upper bound is constructive.
Long output streams are divided into blocks, neighboring \(P\)- and
reversed \(Q\)-blocks are coupled nearly at the larger Shannon entropy,
and the resulting dependence is placed beyond a causal anti-diagonal
that no horizon-\(n\) realized history can cross.
For the crash-test pair,
\[
\boxed{
C_\infty^{\mathrm{causal}}
=
1.5
<
1.728465468745531
\le
C_\infty^{\mathrm{syn}}.
}
\]

Version 5 places these results inside one commutative operator-algebraic
representation.
A finite classical latent law is represented by one probability vector
\(\psi\) and a commutative algebra of multiplication operators.
Stable synchronous identity becomes identification of event projections
across histories.
For each action distribution \(P\), self-information is represented by
a diagonal surprisal operator \(L_P\), whose vector spectral measure is
exactly the surprisal law \(\mu_P\).
Thus entropy, varentropy, higher moments, lattice phase, and the full
characteristic function are different probes of one spectral ledger.

This yields the compact comparison
\[
\boxed{
C_\infty^{\mathrm{causal}}
=
\max_a \langle \psi_a,L_a\psi_a\rangle,
}
\]
for two actions, while
\[
\boxed{
C_\infty^{\mathrm{syn}}
\ge
\frac12
\left[
\langle L_P\rangle+
\langle L_Q\rangle+
d_{\mathrm{Lip}}(P,Q)
\right],
}
\]
where
\[
d_{\mathrm{Lip}}(P,Q)
=
W_1(\mu_P,\mu_Q).
\]
In this precise sense, the causal contract pays asymptotically for the
largest mean surprisal, whereas stable synchronous identity can retain
a cost determined by the full surprisal spectrum.

Finally, we connect this spectral ledger to a scale-orbit calculation
from compactly supported polynomial kernels.
Unitary dilation produces a Mellin spectral measure, and every such
finite-energy orbit measure defines a valid spectral probe of surprisal
mismatch and hence a synchronous lower-bound certificate.
The hard-edged \(k=0\) kernel is singular and reproduces exactly the
\(\min/\max\) scalar identity underlying the Wasserstein converse.
Smoother kernels define a family of softer spectral auditors.

The resulting framework does not identify an ontology and does not
claim that one physical vector underlies reality.
It gives a unified representation of the accounting problem:
one state vector, an algebra acting on it, and a spectral ledger
recording which distinctions the realization contract is required to
preserve.
\end{abstract}


% ============================================================
\section{Introduction: the bill depends on the contract}
% ============================================================

Suppose an observable controlled system is specified by a finite family
of output laws
\[
W=\{P_a:a\in\cA\}.
\]
The system is memoryless in the operational sense that, for every
adaptive nonanticipative action policy,
\[
\Pr(Y_t=y\mid A_{1:t},Y_{1:t-1})
=
P_{A_t}(y).
\tag{1.1}
\]

The observable law does not uniquely determine the latent structure that
realizes it.
A realization may define only the output that is actually requested,
or it may be required to assign definite values to mutually exclusive
alternatives.
If it does the latter, one must still specify whether ``the same''
future event is required to have the same latent identity across
different histories.

The central resource question of this paper is therefore not

\[
\text{``How much randomness does the observable law contain?''}
\]

but

\begin{center}
\fbox{
\begin{minipage}{0.84\textwidth}
\centering
\textbf{What information must an exact realization carry under a
specified counterfactual identity contract?}
\end{minipage}}
\end{center}

The distinction is essential.
The same observable channel will be shown to have different asymptotic
source costs under different latent identity requirements.

The paper has two intertwined layers.

First, an operational counterfactual layer distinguishes a synchronous
stable-event realization from a general history-indexed causal response
tree and computes or bounds their minimum latent Shannon entropies.

Second, an algebraic-spectral layer shows that the principal objects can
be represented by one classical state vector and a commutative algebra
of observables.
The self-information structure of each action is then the spectral
measure of one diagonal operator in one vector state.

The purpose of the second layer is not to replace the first.
It explains why the first layer, the information-spectrum converse, the
kernel scale-reduction calculation, the arithmetic obstruction, and the
causal/synchronous separation belong to one mathematical architecture.


% ============================================================
\section{Priority convention and current literature map}
% ============================================================

This paper uses several standard theories and deliberately avoids
claiming novelty for them.

Minimum-entropy coupling is classical information theory
\cite{Kovacevic2015}.
Information-spectrum converses for minimum-entropy coupling and
functional representation are known \cite{ShkelYadav2023}.
Approximation algorithms for multi-marginal MEC are known
\cite{Li2021}, and the two-marginal profile method gives a
\(\log_2(e)/e\) additive guarantee for the standard greedy coupling
\cite{Compton2023}.
As of 2026, a polynomial-time approximation scheme is also known for a
constant number of marginals \cite{Compton2026PTAS}.
Thus the profile theorem is used below as a clean analytic block-coupling
ingredient, not as a claim about the current algorithmic frontier.

Exact common information and exact channel synthesis provide neighboring
multi-letter viewpoints \cite{KumarLiElGamal,YuTanExact}.
Online entropy-efficient exact sampling and randomness recycling have
advanced substantially, including near-Shannon online samplers and
space--entropy lower bounds \cite{DraperSaad2026,DraperHarrisSaad2026,
DraperSaadSpace2026}.

The relation between scale invariance and Mellin analysis is standard;
a recent systematic tutorial develops the scale-invariant analogue of
shift-invariant signal theory explicitly through Mellin methods
\cite{GuptaSinghAggarwalJoshi2025}.
The use of Hilbert-space distribution embeddings and maximum mean
discrepancy is established \cite{Sriperumbudur2010,Gretton2012}.
Recent work also develops MMD-type two-sample procedures invariant under
general group actions \cite{GiacofciMeynaouiPodgorny2026}.
Zolotarev distances and their Fourier analysis are classical and remain
active tools for higher-order probability comparison
\cite{BobkovLedoux2026}.

Common-random-number implementations with stable event keys are a
particularly close practical neighbor to the distinction between
execution-path identity and event identity made here
\cite{BuffaloPearsonKlein2026}.

The operator-algebraic representation of a classical probability law by
a vector state on a commutative algebra, the spectral theorem, Bochner's
theorem, and the GNS construction are standard mathematics.

Accordingly, throughout the paper:
\begin{quote}
``derived here'' means derived inside the present argument, not claimed
as historically first.
\end{quote}

The specific combination studied here is narrower:
minimum Shannon entropy under all synchronous transversals versus a
history-indexed causal realization; the exact two-action causal
Shannon-floor construction; the finite-horizon residual analysis; and
the use of orbit spectral measures as resolution certificates for this
counterfactual source-cost problem.


% ============================================================
\section{Four realization contracts}
% ============================================================

Let \(\cA\) and \(\cY\) be finite.

\paragraph{C1: chosen-path sampling.}
At time \(t\), the action \(A_t\) is already known.
Only
\[
Y_t\sim P_{A_t}
\]
must be generated.

\paragraph{C2: one-step counterfactual table.}
Before the action is selected, a latent variable fixes
\[
F=(F(a))_{a\in\cA},
\qquad
F(a)\sim P_a.
\tag{3.1}
\]
All action outputs are jointly defined, although only one is exposed.

\paragraph{C3-S: synchronous stable-event table.}
For horizon \(n\), one latent source fixes
\[
F=(F_t(a))_{1\le t\le n,\ a\in\cA}.
\tag{3.2}
\]
The same coordinate \((t,a)\) denotes the same latent event across all
histories.

\paragraph{C3-C: history-indexed causal tree.}
A general deterministic seed fixes
\[
F_t(h_{t-1},a).
\tag{3.3}
\]
The latent event corresponding to the same time/action pair may depend
on the prior history \(h_{t-1}\).

Thus
\[
\boxed{
\text{C3-S is C3-C plus stable cross-history event identity.}
}
\tag{3.4}
\]

The paper asks what additional source entropy that identity can force.


% ============================================================
\section{The one-shot counterfactual bill}
% ============================================================

Define
\[
h_{\max}
:=
\max_{a\in\cA} H(P_a).
\tag{4.1}
\]

Under C2 the minimum latent entropy is
\[
C_1(W)
:=
\min
\left\{
H(F):
\Law(F(a))=P_a,\ \forall a
\right\}.
\tag{4.2}
\]

Hence
\[
\boxed{
C_1(W)
=
\MEC(P_a:a\in\cA).
}
\tag{4.3}
\]

The one-shot surcharge above the largest observable entropy is
\[
r_1(W)
=
C_1(W)-h_{\max}\ge0.
\tag{4.4}
\]

This quantity must not be multiplied by time without proving that the
multi-step realization contract forces independent one-step tables.


% ============================================================
\section{The equal-entropy crash test}
% ============================================================

The concrete test pair used throughout the paper is
\[
P=
\left(
\frac12,\frac14,\frac14
\right),
\qquad
Q=(q,q,1-2q),
\tag{5.1}
\]
where \(q\in(1/3,1/2)\) is the unique solution of
\[
-2q\log_2q-(1-2q)\log_2(1-2q)
=
\frac32.
\tag{5.2}
\]

Numerically,
\[
q
=
0.41003242818198557608\ldots
\tag{5.3}
\]
and
\[
H(P)=H(Q)=1.5.
\tag{5.4}
\]

The probability spectra are not permutations.

A one-shot minimum-entropy coupling has nonzero atom probabilities
\[
q,\qquad
\frac14,\qquad
1-2q,\qquad
\frac12-q,\qquad
2q-\frac34.
\tag{5.5}
\]

Therefore
\[
\boxed{
C_1(P,Q)
=
2.0539185856072597\ldots
}
\tag{5.6}
\]
and
\[
E_1
:=
C_1-1.5
=
0.5539185856072597\ldots.
\tag{5.7}
\]


% ============================================================
\section{Synchronous stable-event blocks}
% ============================================================

For horizon \(n\), let
\[
F=(F_t(a))_{1\le t\le n,\ a\in\cA}.
\]

\begin{definition}[Synchronous policy-safe law]
A law \(\Gamma\) of \(F\) is synchronous policy-safe if
\[
\Law_\Gamma
\left(
F_1(a_1),\ldots,F_n(a_n)
\right)
=
\bigotimes_{t=1}^nP_{a_t}
\tag{6.1}
\]
for every deterministic action word
\[
a^n=(a_1,\ldots,a_n)\in\cA^n.
\]
\end{definition}

\begin{proposition}[Transversal criterion]
A synchronous table law is policy-safe if and only if it realizes the
memoryless controlled channel exactly under every deterministic
nonanticipative policy.
The same holds for randomized policies whose private randomization is
independent of the table.
\end{proposition}

\begin{proof}
Necessity follows from policies that ignore history.

Conversely, fix a deterministic policy and an output string \(y^n\).
Along this output branch the policy determines an action word
\[
a_t=a_t(y^{t-1}).
\]
The branch event is the corresponding transversal event, whose
probability is
\[
\prod_{t=1}^nP_{a_t(y^{t-1})}(y_t)
\]
by (6.1).
\end{proof}

Define
\[
\Csyn_n(W)
:=
\min
\{
H(F):
F\text{ is synchronous policy-safe}
\}.
\tag{6.2}
\]

\begin{proposition}[Subadditivity and prefix lower bound]
For all \(m,n\ge1\),
\[
\Csyn_{m+n}
\le
\Csyn_m+\Csyn_n,
\tag{6.3}
\]
and
\[
\Csyn_{m+n}
\ge
mh_{\max}+\Csyn_n.
\tag{6.4}
\]
Consequently,
\[
\CinfSyn
:=
\lim_{n\to\infty}\frac{\Csyn_n}{n}
=
\inf_{n\ge1}\frac{\Csyn_n}{n}.
\tag{6.5}
\]
\end{proposition}

\begin{proof}
Independent concatenation gives (6.3).

For (6.4), expose an entropy-maximizing action for the first
\(m\) times and retain the complete future \(n\)-step synchronous table.
Conditioned on each realized prefix, the future table remains feasible.
The corresponding projection of the full source therefore has entropy
at least
\[
mh_{\max}+\Csyn_n.
\]
Fekete's lemma yields (6.5).
\end{proof}


% ============================================================
\section{The synchronous observable space}
% ============================================================

Suppose action \(a\) has output alphabet size \(d_a\).
At one time the complete counterfactual state lies in
\[
R=\prod_{a\in\cA}\cY_a.
\]

Let \(V_a\) be the real vector space of functions on \(R\) depending
only on coordinate \(a\), and define
\[
\cO
=
\sum_{a\in\cA}V_a.
\tag{7.1}
\]

\begin{proposition}[Observable-space dimension and block rank]
The one-time observable space has dimension
\[
\boxed{
d_{\mathrm{obs}}
=
1+\sum_{a\in\cA}(d_a-1).
}
\tag{7.2}
\]
The row space generated by all \(n\)-time transversal marginal
constraints is
\[
\cO^{\otimes n}.
\tag{7.3}
\]
Hence
\[
\boxed{
\operatorname{rank}A_n
=
d_{\mathrm{obs}}^n.
}
\tag{7.4}
\]
\end{proposition}

For two ternary actions,
\[
|R|=9,\qquad d_{\mathrm{obs}}=5,
\]
so
\[
\operatorname{rank}A_n=5^n,
\qquad
\dim\ker A_n=9^n-5^n.
\tag{7.5}
\]

The kernel is the linear space in which hidden interactions may live
without changing any permitted transversal marginal.


% ============================================================
\section{Hidden counterfactual synergy}
% ============================================================

Let \(\gamma\) be a one-time coupling of the action marginals on the
complete one-time state \(R\).

Assume there exists bounded nonzero \(h:R\to\R\) such that
\[
\E_\gamma h=0,
\qquad
\E_\gamma[h\mid R(a)]=0
\quad
\forall a.
\tag{8.1}
\]

\begin{proposition}[Hidden-interaction perturbation]
For sufficiently small nonzero \(\varepsilon\),
\[
\Gamma_\varepsilon(r,s)
=
\gamma(r)\gamma(s)
\left(
1+\varepsilon h(r)h(s)
\right)
\tag{8.2}
\]
is a two-step synchronous policy-safe law.
Its one-time marginals are both \(\gamma\), but
\[
I(R_1;R_2)>0,
\]
and therefore
\[
H(R_1,R_2)
<
2H(\gamma).
\tag{8.3}
\]
\end{proposition}

\begin{proof}
For small enough \(|\varepsilon|\), nonnegativity holds.
Every selectable cross-time marginal annihilates the perturbation term
through the conditional-mean assumptions in (8.1).
Thus all transversal laws remain products.
The joint law is nevertheless nonproduct.
\end{proof}


% ============================================================
\subsection{Explicit 23-atom synchronous witness}
% ============================================================

For the crash-test pair, Appendix~\ref{app:witness} gives an explicit
two-step law \(\widetilde G\) with 23 positive atoms satisfying all four
transversal product laws
\[
(P,P),\quad(P,Q),\quad(Q,P),\quad(Q,Q).
\]

Its entropy is
\[
\boxed{
H(\widetilde G)
=
3.878295576158564\ldots.
}
\tag{8.4}
\]

Therefore
\[
\boxed{
\Csyn_2
\le
3.878295576158564
<
2C_1,
}
\tag{8.5}
\]
and
\[
\boxed{
\CinfSyn
\le
1.939147788079282.
}
\tag{8.6}
\]


% ============================================================
\subsection{Net hidden synergy}
% ============================================================

For a synchronous policy-safe law let
\[
R_t=(F_t(a))_{a\in\cA}
\]
be the complete one-time table and define temporal total correlation
\[
\TC(R_{1:n})
=
\sum_{t=1}^nH(R_t)-H(R_{1:n}).
\tag{8.7}
\]

Since \(H(R_t)\ge C_1\),
\[
H(R_{1:n})
=
nC_1
+
\sum_t(H(R_t)-C_1)
-
\TC(R_{1:n}).
\]

Thus
\[
\boxed{
nC_1-\Csyn_n
=
\max_{\Gamma\ \mathrm{policy\text{-}safe}}
\left[
\TC_\Gamma(R_{1:n})
-
\sum_t
\bigl(
H_\Gamma(R_t)-C_1
\bigr)
\right].
}
\tag{8.8}
\]

The saving below \(nC_1\) is therefore temporal hidden dependence minus
the one-time entropy premium required to support that dependence.


% ============================================================
\section{Synchronous information-spectrum geometry}
% ============================================================

For a finite law \(P=(p_i)\), define the surprisal variable
\[
\imath_P(X)
=
-\log_2P(X),
\qquad
X\sim P,
\tag{9.1}
\]
and its law
\[
\mu_P
=
\Law(\imath_P(X)).
\tag{9.2}
\]

\begin{theorem}[Two-action information-spectrum converse]
For every synchronous policy-safe \(n\)-block,
\[
\boxed{
H(F)
\ge
\frac n2
\left[
H(P)+H(Q)+W_1(\mu_P,\mu_Q)
\right].
}
\tag{9.3}
\]
Consequently,
\[
\boxed{
\CinfSyn(P,Q)
\ge
\frac12
\left[
H(P)+H(Q)+W_1(\mu_P,\mu_Q)
\right].
}
\tag{9.4}
\]
\end{theorem}

\begin{proof}
Fix a positive-probability full-table atom
\[
f=(x_1,z_1,\ldots,x_n,z_n)
\]
with probability \(\Gamma(f)\).

For every subset \(S\subseteq[n]\), the same atom lies in the
transversal selecting \(P\) on \(S\) and \(Q\) on \(S^c\).
Therefore
\[
\Gamma(f)
\le
\prod_{t\in S}P(x_t)
\prod_{t\notin S}Q(z_t).
\]
Minimizing over \(S\) separates coordinatewise:
\[
\Gamma(f)
\le
\prod_{t=1}^n
\min\{P(x_t),Q(z_t)\}.
\tag{9.5}
\]
Hence
\[
-\log_2\Gamma(f)
\ge
\sum_{t=1}^n
\max\{
\imath_P(x_t),\imath_Q(z_t)
\}.
\]
Average over the full-table law and use
\[
\max(u,v)
=
\frac12(u+v+|u-v|).
\]
At each time the pair of surprisals is some coupling of
\(\mu_P,\mu_Q\); the minimum expected \(|u-v|\) is their
one-dimensional Wasserstein distance.
\end{proof}

If
\[
H(P)=H(Q)=h,
\]
then
\[
\boxed{
\CinfSyn(P,Q)
\ge
h+\frac12W_1(\mu_P,\mu_Q).
}
\tag{9.6}
\]


% ============================================================
\subsection{Multi-action quantile envelope}
% ============================================================

For each action \(a\), let
\[
U_a=-\log_2P_a(Y_a),
\qquad
Y_a\sim P_a,
\]
with law \(\mu_a\).

Define
\[
\mathcal T(P_a:a\in\cA)
=
\inf_{\eta\in\Pi(\mu_a:a\in\cA)}
\E_\eta\max_{a\in\cA}U_a.
\tag{9.7}
\]

\begin{lemma}[Quantile envelope]
For finite-support laws,
\[
\boxed{
\mathcal T(P_a:a\in\cA)
=
\int_0^1
\max_{a\in\cA}
F_{\mu_a}^{-1}(u)\,\dd u.
}
\tag{9.8}
\]
\end{lemma}

\begin{proof}
For every coupling,
\[
\Pr(\max_aU_a\le t)
\le
\min_aF_{\mu_a}(t).
\]
The comonotone quantile coupling attains equality.
\end{proof}

\begin{theorem}[Multi-action synchronous converse]
\[
\boxed{
\Csyn_n(P_a:a\in\cA)
\ge
n
\int_0^1
\max_{a\in\cA}
F_{\mu_a}^{-1}(u)\,\dd u.
}
\tag{9.9}
\]
\end{theorem}


% ============================================================
\section{The certified synchronous crash-test interval}
% ============================================================

For
\[
P=\left(\frac12,\frac14,\frac14\right)
\]
one has
\[
\mu_P
=
\frac12\delta_1+\frac12\delta_2.
\tag{10.1}
\]

For \(Q=(q,q,1-2q)\),
\[
\mu_Q
=
2q\,\delta_{-\log_2q}
+
(1-2q)\delta_{-\log_2(1-2q)}.
\tag{10.2}
\]

Direct quantile evaluation gives
\[
W_1(\mu_P,\mu_Q)
=
0.4569309374910615\ldots.
\tag{10.3}
\]

Therefore
\[
\boxed{
1.728465468745531
\le
\CinfSyn
\le
1.939147788079282.
}
\tag{10.4}
\]

At depth two,
\[
3.553918585607260
\le
\Csyn_2
\le
3.878295576158564.
\tag{10.5}
\]

The upper endpoint is an explicit witness value, not a proof of
two-step optimality.


% ============================================================
\section{Counterfactual breadth}
% ============================================================

Let
\[
\cG\subseteq\{P,Q\}^n
\]
be only a selected family of action strings required to have the correct
transversal law, and let \(C_n(\cG)\) be the minimum table entropy under
those constraints.

For a full atom define
\[
M_t
=
\max\{\imath_P(X_t),\imath_Q(Z_t)\},
\]
and
\[
\Delta_t
=
|\imath_P(X_t)-\imath_Q(Z_t)|.
\]

Let \(s\in\{P,Q\}^n\) select the locally more surprising action.
For \(g\in\cG\), define
\[
d_\Delta(s,g)
=
\sum_{t:s_t\ne g_t}\Delta_t.
\tag{11.1}
\]

Then
\[
-\log_2\Gamma(f)
\ge
\sum_tM_t-d_\Delta(s,\cG),
\tag{11.2}
\]
where
\[
d_\Delta(s,\cG)
=
\min_{g\in\cG}d_\Delta(s,g).
\]

If
\[
D
=
\max_{x,z}
|\imath_P(x)-\imath_Q(z)|
\tag{11.3}
\]
and
\[
r(\cG)
=
\max_s\min_{g\in\cG}d_H(s,g),
\tag{11.4}
\]
then
\[
\boxed{
C_n(\cG)
\ge
n\mathcal T(P,Q)-Dr(\cG).
}
\tag{11.5}
\]

Thus breadth is governed by covering geometry, not only by the number of
counterfactual strings.


% ============================================================
\section{History-indexed causal trees}
% ============================================================

\begin{definition}[Causal-tree entropy]
Let
\[
\Ccausal_n(W)
\]
be the minimum Shannon entropy of a finite random deterministic
history-indexed response tree whose induced output process realizes the
memoryless controlled channel exactly under every adaptive
nonanticipative policy.
\end{definition}

At horizon one,
\[
\Ccausal_1=C_1.
\tag{12.1}
\]

Every synchronous table is a causal tree via
\[
F_t(h,a)=F_t(a),
\]
so
\[
\boxed{
\Ccausal_n\le\Csyn_n.
}
\tag{12.2}
\]


% ============================================================
\subsection{The action-tree quotient}
% ============================================================

For a deterministic causal seed and action word
\[
a^t=(a_1,\ldots,a_t),
\]
let \(Z_{a^t}\) be the output reached at depth \(t\) by feeding this
action word to the seed.

\begin{proposition}[Action-tree characterization]
A random action-labelled tree is an exact horizon-\(n\) causal
realization if and only if
\[
\boxed{
\Law
\left(
Z_{a_1},
Z_{a_1a_2},
\ldots,
Z_{a_1\cdots a_n}
\right)
=
\bigotimes_{t=1}^nP_{a_t}
}
\tag{12.3}
\]
for every deterministic action word \(a^n\).

Consequently,
\[
\boxed{
\Ccausal_n
=
\min
H\left(
(Z_v)_{v\in\cA^{\le n}}
\right)
}
\tag{12.4}
\]
over all random action trees satisfying (12.3).
\end{proposition}

\begin{proof}
Every history-indexed tree induces such an action tree.

Conversely, for a deterministic adaptive policy and target output
string \(y^n\), that output string determines one action word along the
branch:
\[
a_t=a_t(y^{t-1}).
\]
The target event is exactly the corresponding path event in the action
tree, and (12.3) supplies its product probability.
\end{proof}


% ============================================================
\subsection{Causal prefix recursion and residual entropy}
% ============================================================

\begin{proposition}[Causal prefix lower bound]
For every \(n\ge1\),
\[
\boxed{
\Ccausal_n
\ge
C_1+(n-1)h_{\max}.
}
\tag{12.5}
\]
Moreover,
\[
\Ccausal_{m+n}
\le
\Ccausal_m+\Ccausal_n.
\tag{12.6}
\]
Hence
\[
\CinfCausal
=
\lim_{n\to\infty}\frac{\Ccausal_n}{n}
\tag{12.7}
\]
exists and
\[
h_{\max}
\le
\CinfCausal
\le
\CinfSyn.
\tag{12.8}
\]
\end{proposition}

Assume now that all actions have the same entropy
\[
H(P_a)=h.
\]

Define
\[
E_n
=
\Ccausal_n-nh.
\tag{12.9}
\]

For any fixed deterministic policy \(\pi\), its realized output
trajectory \(Y_\pi^n\) is a deterministic function of the causal seed
\(S_n\), and
\[
H(Y_\pi^n)=nh.
\]
Therefore
\[
\boxed{
H(S_n)
=
nh+H(S_n\mid Y_\pi^n).
}
\tag{12.10}
\]

Minimizing,
\[
\boxed{
E_n
=
\inf_{S_n}
H(S_n\mid Y_\pi^n).
}
\tag{12.11}
\]

The prefix lower bound and block concatenation imply
\[
E_{n+1}\ge E_n,
\qquad
E_{m+n}\le E_m+E_n.
\tag{12.12}
\]

Thus \(E_n\) is nonnegative, nondecreasing, and subadditive, while
\[
\boxed{
\CinfCausal-h
=
\lim_{n\to\infty}\frac{E_n}{n}.
}
\tag{12.13}
\]


% ============================================================
\section{Exact causal Bellman--MEC recursion}
% ============================================================

Let \(\cM_n\) be the set of exact depth-\(n\) action-tree laws.

At a root action \(a\), write
\[
B_a=(Y_a,T_a).
\]
Conditional on \(Y_a=y\), let
\[
K_{a,y}
=
\Law(T_a\mid Y_a=y)
\in\cM_n.
\]

The marginal law of the complete \(a\)-branch is
\[
M_a
=
\bigoplus_yP_a(y)K_{a,y}.
\tag{13.1}
\]

\begin{theorem}[Exact causal Bellman--MEC recursion]
\[
\boxed{
\Ccausal_{n+1}
=
\inf_{\{K_{a,y}\in\cM_n\}}
\MEC(M_a:a\in\cA).
}
\tag{13.2}
\]
\end{theorem}

\begin{proof}
Once the conditional continuation laws are fixed, each root-action
branch has prescribed marginal \(M_a\).
Distinct root actions are mutually exclusive branches of the action
tree and may be coupled arbitrarily.
The least possible entropy of their common latent specification is the
multi-marginal MEC.
Finally minimize over the child laws.
\end{proof}


% ============================================================
\subsection{Shape versus join}
% ============================================================

For equal-entropy actions define
\[
e_a
=
\sum_yP_a(y)
\left[
H(K_{a,y})-\Ccausal_n
\right]\ge0.
\tag{13.3}
\]
Then
\[
H(M_a)
=
h+\Ccausal_n+e_a.
\tag{13.4}
\]

Define
\[
\Delta_{\mathrm{shape}}
=
\max_ae_a
\tag{13.5}
\]
and
\[
\Delta_{\mathrm{join}}
=
\MEC(M_a:a\in\cA)-\max_aH(M_a).
\tag{13.6}
\]

Both are nonnegative, and
\[
\boxed{
E_{n+1}-E_n
=
\inf_K
\left(
\Delta_{\mathrm{shape}}
+
\Delta_{\mathrm{join}}
\right).
}
\tag{13.7}
\]

This identity isolates the finite-horizon causal tension:
continuations can be reshaped to make root branches easier to join, but
departing from optimal child entropy has a cost.


% ============================================================
\section{Finite-depth rigidity}
% ============================================================

For the crash-test pair, the one-step \(3\times3\) transportation
polytope has 18 vertices.
Up to permutation their nonzero spectra fall into five types.
The minimum type is
\[
\lambda_\star
=
\left(
q,\frac14,1-2q,\frac12-q,2q-\frac34
\right).
\tag{14.1}
\]

Appendix~\ref{app:rigidity} verifies that for
\[
\frac38<q<\frac5{12}
\tag{14.2}
\]
this spectrum strictly majorizes each competing vertex spectrum.
By strict Schur concavity, every one-step entropy minimizer therefore
has the same nonzero atom multiset \(\lambda_\star\).

\begin{theorem}[Strict depth-two causal excess]
For the crash-test pair,
\[
\boxed{
\Ccausal_2>h+C_1.
}
\tag{14.3}
\]
Equivalently,
\[
\boxed{
E_2>E_1.
}
\tag{14.4}
\]
\end{theorem}

\begin{proof}
Assume equality.
For root action \(a\), let \(B_a=(Y_a,T_a)\).
Every conditional child law \(K_{a,y}\) is a valid one-step coupling,
so
\[
H(K_{a,y})\ge C_1.
\]
Thus
\[
H(B_a)
=
h+\sum_yP_a(y)H(K_{a,y})
\ge
h+C_1.
\]
Since \(B_a\) is a deterministic function of the full depth-two seed,
global equality forces equality at every stage.
Every \(K_{a,y}\) is therefore a one-step minimizer and has spectrum
\(\lambda_\star\).

Hence
\[
\spec(B_P)
=
\spec(P\otimes\lambda_\star),
\qquad
\spec(B_Q)
=
\spec(Q\otimes\lambda_\star).
\]
Both branch objects have the same entropy as the full seed and are
deterministic images of it, so they must be invertible images on
positive support and therefore have the same atom multiset.
Thus
\[
P\otimes\lambda_\star
\equiv_{\mathrm{perm}}
Q\otimes\lambda_\star.
\]
Power-sum cancellation implies equality of the probability spectra of
\(P\) and \(Q\), contradiction.
\end{proof}

The finite causal excess is therefore real.
The remaining issue is whether it can be made sublinear.


% ============================================================
\section{Low-entropy couplings of long product blocks}
% ============================================================

Let
\[
V_P
=
\Var[-\log_2P(X)],
\qquad
V_Q
=
\Var[-\log_2Q(Y)].
\tag{15.1}
\]

The two-marginal profile method of Compton et al. gives a greedy coupling
\(G(A,B)\) satisfying
\[
H(G(A,B))
\le
H(\operatorname{PROFILE}(A,B))
+
c_0,
\tag{15.2}
\]
where
\[
c_0=\frac{\log_2e}{e}.
\tag{15.3}
\]

For two finite distributions, the profile entropy can be written in
surprisal quantiles as
\[
\boxed{
H(\operatorname{PROFILE}(A,B))
=
\frac12
\left[
H(A)+H(B)+W_1(\mu_A,\mu_B)
\right].
}
\tag{15.4}
\]

\begin{lemma}[Product-block coupling bound]
Let
\[
h_\star=\max\{H(P),H(Q)\},
\qquad
A_0
=
\frac12(\sqrt{V_P}+\sqrt{V_Q}).
\]
For every \(L\ge1\), there exists
\[
J_L\in\Pi(P^{\otimes L},Q^{\otimes L})
\]
such that
\[
\boxed{
H(J_L)
\le
Lh_\star+A_0\sqrt L+c_0.
}
\tag{15.5}
\]
\end{lemma}

\begin{proof}
The product surprisal is a sum of \(L\) i.i.d. surprisals.
By the triangle inequality,
\[
W_1(\mu_{P^L},\mu_{Q^L})
\le
L|H(P)-H(Q)|
+
\sqrt{LV_P}
+
\sqrt{LV_Q}.
\]
Insert this into (15.4) and then use (15.2).
\end{proof}


% ============================================================
\section{The anti-diagonal causal construction}
% ============================================================

Fix
\[
n=mL.
\]
Introduce two length-\(n\) streams
\[
X_1,\ldots,X_n,
\qquad
W_1,\ldots,W_n
\]
divided into \(m\) consecutive blocks
\[
X^{(1)},\ldots,X^{(m)},
\qquad
W^{(1)},\ldots,W^{(m)}
\]
of length \(L\).

Construct the seed from independent components as follows:
\[
X^{(1)}\sim P^L,
\qquad
W^{(m)}\sim Q^L,
\]
and, for every \(j=1,\ldots,m-1\),
\[
(W^{(j)},X^{(j+1)})\sim J_L,
\tag{16.1}
\]
independently across \(j\).

Reverse the \(W\)-stream:
\[
Z_k=W_{n+1-k}.
\tag{16.2}
\]

The causal response rule is:
on the \(r\)-th occurrence of action \(P\), output \(X_r\);
on the \(s\)-th occurrence of action \(Q\), output \(Z_s\).

\begin{lemma}[Anti-diagonal invisibility]
No action history of length at most \(n\) can observe both sides of one
coupled pair
\[
(W^{(j)},X^{(j+1)}).
\]
\end{lemma}

\begin{proof}
To see any coordinate of \(X^{(j+1)}\), the history needs at least
\[
jL+1
\]
occurrences of \(P\).
To see any coordinate of \(W^{(j)}\) through the reversed \(Q\)-stream,
it needs at least
\[
n-jL+1
\]
occurrences of \(Q\).
Seeing both requires at least
\[
(jL+1)+(n-jL+1)=n+2
\]
actions.
\end{proof}

\begin{proposition}[Exactness of the anti-diagonal source]
The construction above is an exact horizon-\(n\) history-indexed causal
realization of the memoryless channel \(\{P,Q\}\).
\end{proposition}

\begin{proof}
Fix a deterministic action word.
By the anti-diagonal lemma, the realized path exposes at most one
marginal side of each coupled component.
That exposed side has exactly product law \(P^L\) or \(Q^L\), and the
different components are independent.
Hence the complete realized path has the required product law.

For an adaptive policy and a target output string, the output string
determines one action word along the branch, so the same argument
applies.
\end{proof}

The source entropy is
\[
H(S_n)
=
LH(P)+LH(Q)+(m-1)H(J_L).
\tag{16.3}
\]

Therefore
\[
\boxed{
H(S_n)
\le
nh_\star
+
Lh_{\min}
+
(m-1)(A_0\sqrt L+c_0),
}
\tag{16.4}
\]
where
\[
h_{\min}=\min\{H(P),H(Q)\}.
\]


% ============================================================
\section{The causal Shannon-floor theorem}
% ============================================================

\begin{theorem}[Causal Shannon floor for two actions]
For arbitrary finite probability distributions \(P,Q\),
\[
\boxed{
\CinfCausal(P,Q)
=
\max\{H(P),H(Q)\}.
}
\tag{17.1}
\]
\end{theorem}

\begin{proof}
Let
\[
h_\star=\max\{H(P),H(Q)\}.
\]

The fixed policy that always selects an entropy-maximizing action
produces an \(n\)-step output trajectory of entropy \(nh_\star\).
Since this trajectory is a deterministic projection of the causal seed,
\[
\Ccausal_n\ge nh_\star.
\]

For the upper bound, use the anti-diagonal construction with
\(n=mL\).
By (16.4),
\[
\frac{\Ccausal_{mL}}{mL}
\le
h_\star
+
\frac{h_{\min}}m
+
\frac{A_0}{\sqrt L}
+
\frac{c_0}{L}.
\]
Let \(m,L\to\infty\).
\end{proof}

\begin{corollary}[Sublinear finite-horizon residual]
For fixed finite \(P,Q\),
\[
\boxed{
\Ccausal_n
\le
nh_\star+O(n^{2/3}).
}
\tag{17.2}
\]
For equal-entropy actions,
\[
\boxed{
E_n=O(n^{2/3}).
}
\tag{17.3}
\]
\end{corollary}

\begin{proof}
Take \(L\asymp n^{2/3}\) and \(m\asymp n^{1/3}\) in
(16.4).
A terminal incomplete block costs only \(O(L)\).
\end{proof}


% ============================================================
\section{Strict asymptotic separation for the crash test}
% ============================================================

Since
\[
H(P)=H(Q)=1.5,
\]
the causal Shannon-floor theorem gives
\[
\boxed{
\CinfCausal=1.5.
}
\tag{18.1}
\]

Together with the synchronous interval,
\[
\boxed{
1.5
=
\CinfCausal
<
1.728465468745531
\le
\CinfSyn
\le
1.939147788079282
<
C_1.
}
\tag{18.2}
\]

Hence
\[
\boxed{
\CinfSyn-\CinfCausal
\ge
0.228465468745531.
}
\tag{18.3}
\]

The gap is not a cost of adaptive sampling itself.
It is certified under the stronger requirement that the same
time/action event retain one common latent identity across mutually
exclusive histories.


% ============================================================
\section{Finite non-attainment and asymptotic attainment}
% ============================================================

For the crash test,
\[
E_1>0.
\]
Since \(E_n\) is nondecreasing,
\[
E_n\ge E_1
\qquad
\forall n.
\tag{19.1}
\]

Thus
\[
\boxed{
\Ccausal_n>1.5n
\qquad
\text{for every finite }n.
}
\tag{19.2}
\]

Nevertheless,
\[
E_n=O(n^{2/3}),
\]
so
\[
\frac{E_n}{n}\to0.
\tag{19.3}
\]

The Shannon floor is therefore an asymptotic rate, not a finite-horizon
exact equality.


% ============================================================
\section{The one-vector classical representation}
\label{sec:onevector}
% ============================================================

We now recast the finite latent model algebraically.

Let \(\Omega=\supp\Gamma\) be the positive support of a finite latent
law \(\Gamma\).
Define
\[
\cH_\Gamma=\ell^2(\Omega)
\]
and
\[
\boxed{
\psi_\Gamma(\omega)=\sqrt{\Gamma(\omega)}.
}
\tag{20.1}
\]

Every real random variable \(f:\Omega\to\R\) defines the multiplication
operator
\[
M_f\xi(\omega)
=
f(\omega)\xi(\omega).
\tag{20.2}
\]

Then
\[
\boxed{
\E_\Gamma f
=
\ip{\psi_\Gamma}{M_f\psi_\Gamma}.
}
\tag{20.3}
\]

For an event \(E\subseteq\Omega\),
\[
\Pi_E=M_{\1_E}
\]
is an orthogonal projection and
\[
\boxed{
\Gamma(E)
=
\ip{\psi_\Gamma}{\Pi_E\psi_\Gamma}.
}
\tag{20.4}
\]

Thus every finite classical source may be represented as
\[
\boxed{
\text{one normalized vector }\psi_\Gamma
+
\text{the commutative algebra }
\cA_\Gamma\simeq\ell^\infty(\Omega).
}
\tag{20.5}
\]

Because \(\Gamma\) has positive support, \(\psi_\Gamma\) is cyclic for
the multiplication algebra:
\[
\overline{\cA_\Gamma\psi_\Gamma}
=
\cH_\Gamma.
\tag{20.6}
\]

\begin{remark}[The priced entropy remains Shannon entropy]
This is a Hilbert-space representation of classical probability.
The algebra is commutative and no quantum probability is introduced.

The resource priced in this paper is
\[
H(\Gamma)
=
-\sum_\omega
|\psi_\Gamma(\omega)|^2
\log_2|\psi_\Gamma(\omega)|^2.
\tag{20.7}
\]
It is not the von Neumann entropy of the rank-one operator
\(
|\psi_\Gamma\rangle\langle\psi_\Gamma|
\).
\end{remark}


% ============================================================
\section{Stable identity as operator identification}
% ============================================================

For a synchronous table define the event projections
\[
E_{t,a,y}
=
\1\{F_t(a)=y\}
\in\cA_\Gamma.
\tag{21.1}
\]

The transversal condition becomes
\[
\boxed{
\ip{\psi_\Gamma}{
\left(
\prod_{t=1}^n E_{t,a_t,y_t}
\right)
\psi_\Gamma
}
=
\prod_{t=1}^nP_{a_t}(y_t).
}
\tag{21.2}
\]

The defining identity requirement of C3-S is that
\[
E_{t,a,y}
\]
is literally the same projection for all histories that might precede
time \(t\).

For C3-C the corresponding projection is
\[
E_{t,h,a,y},
\tag{21.3}
\]
and different histories may induce different projections.

Hence
\[
\boxed{
\begin{aligned}
\text{C3-S:}\quad&
E_{t,h,a,y}\equiv E_{t,a,y}
\quad\forall h,
\\
\text{C3-C:}\quad&
E_{t,h,a,y}
\text{ may depend on }h.
\end{aligned}
}
\tag{21.4}
\]

Thus stable synchronous identity is an algebraic identification of
observables across mutually exclusive histories.


% ============================================================
\section{The surprisal operator and spectral ledger}
% ============================================================

Let
\[
P=(p_1,\ldots,p_N)
\]
have positive support.
Define
\[
\psi_P
=
(\sqrt{p_1},\ldots,\sqrt{p_N})^\top
\in\C^N
\tag{22.1}
\]
and the diagonal surprisal operator
\[
\boxed{
L_P
=
\operatorname{diag}
(-\log_2p_1,\ldots,-\log_2p_N).
}
\tag{22.2}
\]

The vector spectral measure of \(L_P\) in \(\psi_P\) is exactly
\[
\mu_P=\Law[-\log_2P(X)].
\tag{22.3}
\]

Indeed, for every Borel set \(B\),
\[
\boxed{
\mu_P(B)
=
\ip{\psi_P}{E_{L_P}(B)\psi_P}.
}
\tag{22.4}
\]

Therefore
\[
\boxed{
H(P)
=
\ip{\psi_P}{L_P\psi_P},
}
\tag{22.5}
\]
and
\[
\boxed{
V_P
=
\ip{\psi_P}{
(L_P-H(P)I)^2\psi_P
}.
}
\tag{22.6}
\]

More generally,
\[
\E[\imath_P(X)^m]
=
\ip{\psi_P}{L_P^m\psi_P}.
\tag{22.7}
\]

The full characteristic function is
\[
\boxed{
\Phi_P(t)
=
\ip{\psi_P}{e^{itL_P}\psi_P}
=
\int e^{itu}\,\dd\mu_P(u).
}
\tag{22.8}
\]

If
\[
Z_P(s)
=
\sum_jp_j^s,
\tag{22.9}
\]
then
\[
\boxed{
\Phi_P(t)
=
Z_P
\left(
1-\frac{it}{\ln2}
\right).
}
\tag{22.10}
\]

We call \(\Phi_P\), equivalently \(\mu_P\), the
\emph{surprisal spectral ledger}.

It is complete for finite probability spectra:
if \(u=-\log_2p\) is an atom of \(\mu_P\) with mass \(m_u\), then the
multiplicity of the original probability \(p\) is
\[
\boxed{
N_p=m_u2^u.
}
\tag{22.11}
\]
Thus \(\mu_P\) determines the nonzero probability multiset of \(P\) up
to permutation.


% ============================================================
\section{Lipschitz distance between spectral states}
% ============================================================

Define
\[
\boxed{
d_{\mathrm{Lip}}(P,Q)
=
\sup_{\operatorname{Lip}(f)\le1}
\left|
\ip{\psi_P}{f(L_P)\psi_P}
-
\ip{\psi_Q}{f(L_Q)\psi_Q}
\right|.
}
\tag{23.1}
\]

By the spectral theorem and Kantorovich--Rubinstein duality,
\[
\boxed{
d_{\mathrm{Lip}}(P,Q)
=
W_1(\mu_P,\mu_Q).
}
\tag{23.2}
\]

The two principal asymptotic theorems can therefore be written in one
language.

For two actions,
\[
\boxed{
\CinfCausal(P,Q)
=
\max
\{
\ip{\psi_P}{L_P\psi_P},
\ip{\psi_Q}{L_Q\psi_Q}
\}.
}
\tag{23.3}
\]

The synchronous converse gives
\[
\boxed{
\CinfSyn(P,Q)
\ge
\frac12
\left[
\ip{\psi_P}{L_P\psi_P}
+
\ip{\psi_Q}{L_Q\psi_Q}
+
d_{\mathrm{Lip}}(P,Q)
\right].
}
\tag{23.4}
\]

For equal entropy,
\[
\boxed{
\CinfCausal=h,
\qquad
\CinfSyn
\ge
h+\frac12d_{\mathrm{Lip}}(P,Q).
}
\tag{23.5}
\]

This is the central algebraic reading of Version 5:

\begin{center}
\fbox{
\begin{minipage}{0.84\textwidth}
For a finite two-action channel, the history-indexed causal rate is
determined by the largest first spectral moment of surprisal.

The stable synchronous identity contract can remain sensitive to the
shape of the complete surprisal spectral measure.
\end{minipage}}
\end{center}


% ============================================================
\section{Moment blindness as incomplete functional calculus}
% ============================================================

Matching the first \(R\) surprisal moments means
\[
\ip{\psi_P}{L_P^m\psi_P}
=
\ip{\psi_Q}{L_Q^m\psi_Q},
\qquad
m=0,\ldots,R.
\tag{24.1}
\]

Thus the two spectral states agree on the finite polynomial family
\[
I,L,\ldots,L^R.
\]

They need not agree on the full functional calculus.
For example,
\[
f(L)=e^{i\pi L}
\tag{24.2}
\]
can detect a parity distinction invisible to all those moments.

Hence
\[
\boxed{
\text{finite moment information}
=
\text{finite polynomial accounting},
}
\tag{24.3}
\]
while
\[
\boxed{
\text{the full spectral ledger}
=
\text{complete functional calculus}.
}
\tag{24.4}
\]


% ============================================================
\section{Block surprisal resolution}
% ============================================================

Define
\[
\boxed{
\mathcal R_L(P,Q)
=
\frac1L
W_1(\mu_P^{*L},\mu_Q^{*L}).
}
\tag{25.1}
\]

This is the per-symbol mismatch after aggregating \(L\) independent
surprisals.

\begin{theorem}[Block-resolution trichotomy]
Let \(P,Q\) be finite distributions.

If
\[
H(P)\ne H(Q),
\]
then
\[
\boxed{
\mathcal R_L(P,Q)
\to
|H(P)-H(Q)|.
}
\tag{25.2}
\]

If
\[
H(P)=H(Q)
\]
but
\[
V_P\ne V_Q,
\]
then
\[
\boxed{
\sqrt L\,\mathcal R_L(P,Q)
\to
\sqrt{\frac2\pi}
\left|
\sqrt{V_P}-\sqrt{V_Q}
\right|.
}
\tag{25.3}
\]

If
\[
H(P)=H(Q)
\quad\text{and}\quad
V_P=V_Q,
\]
then
\[
\boxed{
\mathcal R_L(P,Q)=O(L^{-1}).
}
\tag{25.4}
\]
\end{theorem}

\begin{proof}
For the first case,
\[
W_1(\mu_P^{*L},\mu_Q^{*L})
\ge
L|H(P)-H(Q)|
\]
by the difference of means, while centered fluctuations contribute only
\(O(\sqrt L)\), giving (25.2).

For equal entropy, center both sums and divide by \(\sqrt L\).
Finite support implies convergence in \(W_1\) to the corresponding
Gaussian laws.
If the limiting variances differ, continuity of \(W_1\) and the monotone
Gaussian coupling give
\[
W_1(N(0,V_P),N(0,V_Q))
=
\sqrt{\frac2\pi}
|\sqrt{V_P}-\sqrt{V_Q}|,
\]
which yields (25.3).

If the variances also agree, a quantitative \(W_1\) central limit
estimate gives normalized error \(O(L^{-1/2})\) for each law.
After restoring the \(\sqrt L\) scale, their unnormalized
Wasserstein distance is \(O(1)\), proving (25.4).
\end{proof}


% ============================================================
\subsection{Crash-test block coefficient}
% ============================================================

For the crash-test pair,
\[
V_P=0.25,
\qquad
V_Q=0.208347310239073\ldots.
\]
Hence
\[
\boxed{
W_1(\mu_P^{*L},\mu_Q^{*L})
=
c_{\mathrm{crash}}\sqrt L+o(\sqrt L),
}
\tag{25.5}
\]
with
\[
c_{\mathrm{crash}}
=
\sqrt{\frac2\pi}
\left(
\frac12-\sqrt{V_Q}
\right)
=
0.0347472540518\ldots.
\tag{25.6}
\]

The small-\(L\) sequence need not be monotone; the statement is
asymptotic.


% ============================================================
\subsection{Matched varentropy improves the causal residual}
% ============================================================

\begin{theorem}[Matched-varentropy causal refinement]
If
\[
H(P)=H(Q)=h,
\qquad
V_P=V_Q,
\]
then
\[
\boxed{
\Ccausal_n
\le
nh+O(\sqrt n).
}
\tag{25.7}
\]
\end{theorem}

\begin{proof}
The third case of the block-resolution theorem gives
\[
W_1(\mu_P^{*L},\mu_Q^{*L})=O(1).
\]
The profile theorem therefore supplies block couplings with
\[
H(J_L)\le Lh+O(1).
\]
The anti-diagonal construction yields
\[
E_n
\le
O(L)+O(n/L).
\]
Choose \(L\asymp\sqrt n\).
\end{proof}


% ============================================================
\section{An arithmetic barrier}
% ============================================================

Higher moment matching does not by itself force a resolution exponent
larger than one in raw \(W_1\).

\begin{lemma}[Lattice-phase lower bound]
Suppose
\[
\supp\mu\subseteq a+h\mathbb Z,
\qquad
\supp\nu\subseteq b+h\mathbb Z.
\]
Then
\[
\boxed{
W_1(\mu^{*L},\nu^{*L})
\ge
\operatorname{dist}
(L(a-b),h\mathbb Z).
}
\tag{26.1}
\]
\end{lemma}

\begin{proof}
The convolution supports lie in
\[
La+h\mathbb Z
\quad\text{and}\quad
Lb+h\mathbb Z,
\]
so every coupled pair is separated by at least the displayed coset
distance.
\end{proof}


% ============================================================
\subsection{Dyadic parity ladder}
% ============================================================

Fix \(R\ge2\).
For \(j=0,\ldots,R+1\), set
\[
p_j=2^{-(R+j)}.
\]

Define \(P_R\) by including, for even \(j\),
\[
2^j\binom{R+1}{j}
\]
atoms of probability \(p_j\), and define \(Q_R\) analogously using odd
\(j\).

Their surprisal laws are
\[
\mu_{P_R}
=
2^{-R}
\sum_{\substack{0\le j\le R+1\\j\ \mathrm{even}}}
\binom{R+1}{j}
\delta_{R+j},
\tag{26.2}
\]
and
\[
\mu_{Q_R}
=
2^{-R}
\sum_{\substack{0\le j\le R+1\\j\ \mathrm{odd}}}
\binom{R+1}{j}
\delta_{R+j}.
\tag{26.3}
\]

\begin{theorem}[Arbitrarily high finite-moment camouflage]
For every \(R\ge2\),
\[
\boxed{
\int u^m\,\dd\mu_{P_R}(u)
=
\int u^m\,\dd\mu_{Q_R}(u),
\qquad
m=0,\ldots,R.
}
\tag{26.4}
\]

Nevertheless,
\[
\boxed{
W_1(\mu_{P_R},\mu_{Q_R})=1.
}
\tag{26.5}
\]

Moreover, for every odd \(L\),
\[
\boxed{
W_1(\mu_{P_R}^{*L},\mu_{Q_R}^{*L})\ge1.
}
\tag{26.6}
\]
\end{theorem}

\begin{proof}
Subtracting the two laws gives the signed measure
\[
2^{-R}
\sum_{j=0}^{R+1}
(-1)^j
\binom{R+1}{j}
\delta_{R+j}.
\]
Its integral against every polynomial of degree at most \(R\) vanishes
because this is an \((R+1)\)-st finite difference.

For the Wasserstein upper bound, identify the weights with the Hamming
weight of a uniformly random \((R+1)\)-bit string conditioned on even or
odd parity.
Flipping one fixed bit maps the even class bijectively to the odd class
and changes Hamming weight by exactly one.
Thus a coupling with distance one exists.

The two supports lie on opposite integer parity classes, so every
coupling has distance at least one.
This proves (26.5).

For odd \(L\), the convolution supports remain on opposite parity
classes, giving (26.6).
\end{proof}

Their common entropy and varentropy are
\[
\boxed{
H(P_R)=H(Q_R)=\frac{3R+1}{2},
}
\tag{26.7}
\]
and
\[
\boxed{
V_{P_R}=V_{Q_R}=\frac{R+1}{4}.
}
\tag{26.8}
\]

Thus arbitrarily many matching moments can coexist with an
\(O(1)\) lattice-phase obstruction to block Wasserstein distance.


% ============================================================
\section{Finite-moment blindness of the synchronous tax}
% ============================================================

For every \(R\ge2\), the causal theorem gives
\[
\CinfCausal(P_R,Q_R)
=
\frac{3R+1}{2}.
\tag{27.1}
\]

Since
\[
W_1(\mu_{P_R},\mu_{Q_R})=1,
\]
the synchronous converse gives
\[
\boxed{
\CinfSyn(P_R,Q_R)
\ge
\frac{3R+1}{2}+\frac12.
}
\tag{27.2}
\]

Hence
\[
\boxed{
\CinfSyn(P_R,Q_R)
-
\CinfCausal(P_R,Q_R)
\ge
\frac12
}
\tag{27.3}
\]
while the first \(R\) surprisal moments agree.

Therefore no fixed finite truncation of the moment sequence determines
the presence or absence of a positive synchronous stable-identity
surcharge.


% ============================================================
\section{Higher-order metrics and the limit of raw Wasserstein}
% ============================================================

A tempting rule would be
\[
\text{first unmatched cumulant of order }r
\Longrightarrow
\mathcal R_L
\asymp
L^{-(r-1)/2}.
\tag{28.1}
\]

The parity ladder disproves this as a universal statement for raw
\(W_1\).
Finite moment matching controls the low-frequency Taylor expansion of
the characteristic function near the origin, while \(W_1\) can also
detect nonzero-frequency lattice phase.

A classical framework for moment-sensitive comparison is provided by
Zolotarev ideal metrics.
For suitable \(s>0\), these metrics satisfy convolution and scaling
properties of the form
\[
\zeta_s(X+Z,Y+Z)\le\zeta_s(X,Y)
\tag{28.2}
\]
for common independent \(Z\), and
\[
\zeta_s(cX,cY)
=
|c|^s\zeta_s(X,Y).
\tag{28.3}
\]

They therefore provide a genuine higher-order smoothing hierarchy once
the necessary lower moments agree.

The open issue for the present project is operational:
find a higher-order resolution functional that both enjoys such ideal
smoothing properties and controls a minimum-entropy coupling or
counterfactual source cost strongly enough to replace the present
\(W_1\) interface.


% ============================================================
\section{One-parameter orbits and spectral measures}
% ============================================================

Let
\[
U_t=e^{itA},
\qquad
t\in\R,
\]
be a strongly continuous one-parameter unitary group on a Hilbert space
\(\cH\), with self-adjoint generator \(A\).
For a normalized vector \(\psi\), define
\[
\boxed{
R_{\psi,A}(t)
=
\ip{\psi}{e^{itA}\psi}.
}
\tag{29.1}
\]

By the spectral theorem there exists a unique probability measure
\[
\nu_{\psi,A}
\]
such that
\[
\boxed{
R_{\psi,A}(t)
=
\int_\R e^{it\xi}\,\dd\nu_{\psi,A}(\xi).
}
\tag{29.2}
\]

\begin{theorem}[One-orbit reduction]
For every \(s,t\in\R\),
\[
\boxed{
\ip{U_s\psi}{U_t\psi}
=
R_{\psi,A}(t-s).
}
\tag{29.3}
\]
Thus pairwise comparison along one unitary orbit depends only on the
relative group displacement.
\end{theorem}

This is the abstract form of the scale reduction used below.


% ============================================================
\section{The dilation model and the REAL kernel family}
% ============================================================

For \(K\in L^2(\R^d)\), define
\[
(U_tK)(x)
=
e^{-dt/2}K(e^{-t}x).
\tag{30.1}
\]
This is a unitary dilation group.

Its infinitesimal generator is, up to the conventional factor \(-i\),
\[
\boxed{
\mathcal D
=
x\cdot\nabla+\frac d2.
}
\tag{30.2}
\]

Let
\[
K_a(x)=a^{-d}K(x/a).
\tag{30.3}
\]
Then
\[
K_a
=
a^{-d/2}U_{\ln a}K.
\tag{30.4}
\]

Hence
\[
\boxed{
\ip{K_a}{K_b}
=
(ab)^{-d/2}
\|K\|_2^2
R_K\left(\ln\frac ba\right),
}
\tag{30.5}
\]
where
\[
R_K(t)
=
\ip{K/\|K\|_2}{U_tK/\|K\|_2}.
\]

For \(0<a\le b\), setting
\[
r=\frac ab
\]
gives the scale-reduction form
\[
\boxed{
\ip{K_a}{K_b}
=
b^{-d}H_K(r),
}
\tag{30.6}
\]
with
\[
H_K(r)
=
r^{-d/2}\|K\|_2^2R_K(-\ln r).
\tag{30.7}
\]

Thus the two absolute scales reduce exactly to one dimensionless ratio.


% ============================================================
\subsection{Compactly supported polynomial kernels}
% ============================================================

In one dimension consider
\[
\rho_k(x)
=
\frac{k+1}{2}(1-|x|)^k_+,
\qquad
k=0,1,2,\ldots.
\tag{30.8}
\]

For
\[
\rho_{k,a}(x)
=
a^{-1}\rho_k(x/a),
\]
define
\[
F_k(a,b)
=
\int_\R\rho_{k,a}(x)\rho_{k,b}(x)\,\dd x.
\tag{30.9}
\]

For \(0<a\le b\),
\[
\boxed{
F_k(a,b)
=
\frac1bH_k(a/b),
}
\tag{30.10}
\]
where
\[
H_k(r)
=
2C_k^2
\int_0^1
(1-u)^k(1-ru)^k\,\dd u,
\qquad
C_k=\frac{k+1}{2}.
\tag{30.11}
\]

Equivalently,
\[
H_k(r)
=
2C_k^2
\sum_{j=0}^{k}
(-1)^j
\binom{k}{j}
B(j+1,k+1)r^j,
\tag{30.12}
\]
and
\[
H_k(r)
=
\frac{k+1}{2}
\,{}_2F_1(-k,1;k+2;r),
\tag{30.13}
\]
where the integral representation (30.11) is the primary form
used below.

For \(m\le k\),
\[
\boxed{
(-1)^mH_k^{(m)}(r)\ge0.
}
\tag{30.14}
\]

The endpoint values are
\[
H_k(1)
=
\frac{(k+1)^2}{2(2k+1)},
\tag{30.15}
\]
and
\[
H_k'(1)
=
-\frac{(k+1)^2}{4(2k+1)}.
\tag{30.16}
\]



% ============================================================
\section{Explicit Mellin spectrum of the polynomial kernels}
% ============================================================

The dilation group is diagonalized by the Mellin transform on each
half-line.
For the even kernel \(\rho_k\), the normalized vector spectral measure
has density

\[
\boxed{
w_k(\xi)
=
\frac{2k+1}{2\pi}
\left|
B\left(\frac12-i\xi,k+1\right)
\right|^2.
}
\tag{31.1}
\]

For integer \(k\),
\[
\boxed{
w_k(\xi)
=
\frac{2k+1}{2\pi}
\frac{(k!)^2}
{\displaystyle
\prod_{j=0}^{k}
\left[
\xi^2+\left(j+\frac12\right)^2
\right]}.
}
\tag{31.2}
\]

It is normalized:
\[
\int_\R w_k(\xi)\,\dd\xi=1,
\tag{31.3}
\]
and
\[
\boxed{
R_k(t)
=
\int_\R e^{it\xi}w_k(\xi)\,\dd\xi.
}
\tag{31.4}
\]

\begin{proof}[Derivation]
On \(x>0\),
\[
\rho_k(x)=C_k(1-x)^k\1_{(0,1)}(x).
\]
The unitary Mellin transform gives
\[
\widehat\rho_k(\xi)
=
\frac{C_k}{\sqrt{2\pi}}
B\left(\frac12-i\xi,k+1\right).
\]
Because the full kernel is even, the positive and negative half-line
contributions are identical.
Dividing the doubled Mellin energy by
\[
\|\rho_k\|_2^2
=
\frac{(k+1)^2}{2(2k+1)}
\]
gives (31.1).
The product formula follows from
\[
\frac{\Gamma(\frac12-i\xi)}
{\Gamma(k+\frac32-i\xi)}
=
\prod_{j=0}^{k}
\frac1{j+\frac12-i\xi}.
\]
\end{proof}

The tail is
\[
\boxed{
w_k(\xi)\asymp|\xi|^{-2k-2}.
}
\tag{31.5}
\]

Therefore
\[
\boxed{
\int|\xi|^\alpha w_k(\xi)\,\dd\xi<\infty
\iff
\alpha<2k+1.
}
\tag{31.6}
\]

The kernel edge regularity is thus recorded directly in the Mellin
spectral tail.


% ============================================================
\section{Hard and smooth orbit-resolution classes}
% ============================================================

For \(k=0\),
\[
w_0(\xi)
=
\frac1{2\pi(\xi^2+1/4)},
\tag{32.1}
\]
the Cauchy law of scale \(1/2\).
Hence
\[
\boxed{
R_0(t)=e^{-|t|/2}.
}
\tag{32.2}
\]

The correlation profile has a linear cusp at the origin and no finite
second spectral moment.

For \(k\ge1\), the second spectral moment is finite and equals
\[
\boxed{
\mathcal E_k
:=
\int\xi^2w_k(\xi)\,\dd\xi
=
\frac{2k+1}{4(2k-1)}.
}
\tag{32.3}
\]

Equivalently,
\[
\mathcal E_k
=
\frac{
\|\left(x\partial_x+\frac12\right)\rho_k\|_2^2
}{
\|\rho_k\|_2^2
}.
\tag{32.4}
\]

Thus
\[
\boxed{
R_k(t)
=
1-\frac{\mathcal E_k}{2}t^2+o(t^2),
\qquad
k\ge1.
}
\tag{32.5}
\]

The hard edge and smooth edge therefore form distinct local resolution
classes:
\[
\boxed{
\begin{array}{ccl}
k=0
&:&
1-R_0(t)\asymp|t|,
\\[1mm]
k\ge1
&:&
1-R_k(t)\asymp t^2.
\end{array}
}
\tag{32.6}
\]


% ============================================================
\section{Reciprocal probability scaling}
% ============================================================

Define the normalized reciprocal-scale overlap
\[
\boxed{
\mathcal A_k(p,q)
=
\frac{F_k(1/p,1/q)}{H_k(1)}.
}
\tag{33.1}
\]

Let
\[
m=\min(p,q),
\qquad
M=\max(p,q).
\]
Then
\[
\boxed{
\mathcal A_k(p,q)
=
m
\frac{H_k(m/M)}{H_k(1)}.
}
\tag{33.2}
\]

Equivalently,
\[
\boxed{
\mathcal A_k(p,q)
=
\sqrt{pq}\,
R_k(|\ln(p/q)|).
}
\tag{33.3}
\]

Introduce surprisal coordinates
\[
u=-\log_2p,
\qquad
v=-\log_2q
\]
and define
\[
c_k(u,v)
=
-\log_2\mathcal A_k(p,q).
\tag{33.4}
\]

Then
\[
\boxed{
c_k(u,v)
=
\frac{u+v}{2}
-
\log_2
R_k((\ln2)|u-v|).
}
\tag{33.5}
\]

For \(k=0\),
\[
\boxed{
\mathcal A_0(p,q)=\min(p,q),
}
\tag{33.6}
\]
and therefore
\[
\boxed{
c_0(u,v)=\max(u,v).
}
\tag{33.7}
\]

This is the exact scalar identity underlying the synchronous
information-spectrum converse.

For \(k\ge1\), using (32.5),
\[
\boxed{
c_k(u,v)
=
\frac{u+v}{2}
+
\frac{\ln2}{2}\mathcal E_k(u-v)^2
+
o((u-v)^2).
}
\tag{33.8}
\]

Explicitly,
\[
\boxed{
c_k(u,v)
=
\frac{u+v}{2}
+
\frac{\ln2(2k+1)}{8(2k-1)}
(u-v)^2
+
o((u-v)^2).
}
\tag{33.9}
\]

Thus the local information-resolution curvature is proportional to the
dilation-generator energy.


% ============================================================
\section{Hard-to-soft overlap ladder}
% ============================================================

Since \(H_k\) is decreasing on \((0,1]\),
\[
H_k(r)\ge H_k(1).
\]
Therefore
\[
\mathcal A_k(p,q)\ge\min(p,q).
\tag{34.1}
\]

Cauchy--Schwarz gives
\[
\mathcal A_k(p,q)\le\sqrt{pq}.
\tag{34.2}
\]

Hence
\[
\boxed{
\min(p,q)
\le
\mathcal A_k(p,q)
\le
\sqrt{pq},
}
\tag{34.3}
\]
and
\[
\boxed{
\frac{u+v}{2}
\le
c_k(u,v)
\le
\max(u,v).
}
\tag{34.4}
\]

Therefore the higher-\(k\) overlap costs cannot strengthen the direct
\(k=0\) synchronous converse.
They form softer, smoother resolution probes.

The large-\(k\) limit is explicit.
From (30.11),
\[
\boxed{
\frac{H_k(r)}{H_k(1)}
\longrightarrow
\frac{2}{1+r}
}
\tag{34.5}
\]
for fixed \(0\le r\le1\).

Consequently,
\[
\boxed{
\mathcal A_k(p,q)
\longrightarrow
\frac{2pq}{p+q},
}
\tag{34.6}
\]
the harmonic mean of \(p\) and \(q\).

In surprisal coordinates,
\[
\boxed{
c_\infty(u,v)
=
\frac{u+v}{2}
+
\log_2
\cosh
\left(
\frac{\ln2}{2}(u-v)
\right).
}
\tag{34.7}
\]

At the orbit level,
\[
\boxed{
R_k(t)\longrightarrow\operatorname{sech}(t/2),
}
\tag{34.8}
\]
whose spectral density under the present Fourier convention is
\[
\boxed{
w_\infty(\xi)=\operatorname{sech}(\pi\xi).
}
\tag{34.9}
\]


% ============================================================
\section{Soft synchronous transport costs}
% ============================================================

For every full synchronous atom,
\[
\Gamma(f)
\le
\prod_t\min\{P(x_t),Q(z_t)\}.
\]
Since
\[
\min(p,q)\le\mathcal A_k(p,q),
\]
we also have the weaker but valid pointwise bound
\[
-\log_2\Gamma(f)
\ge
\sum_t c_k(
\imath_P(x_t),
\imath_Q(z_t)
).
\]

Define
\[
\boxed{
\mathcal T_k(P,Q)
=
\inf_{\pi\in\Pi(\mu_P,\mu_Q)}
\E_\pi c_k(U,V).
}
\tag{35.1}
\]

Then
\[
\boxed{
\Csyn_n(P,Q)
\ge
n\mathcal T_k(P,Q).
}
\tag{35.2}
\]

For \(k=0\),
\[
\mathcal T_0(P,Q)
=
\frac12
\left[
H(P)+H(Q)+W_1(\mu_P,\mu_Q)
\right].
\tag{35.3}
\]

The \(k\ge1\) family is generally weaker as a lower bound, but smoother
as a function of local surprisal mismatch.
It is useful as a bridge to finite-energy orbit probes rather than as an
attempt to beat the hard \(k=0\) converse.


% ============================================================
\section{Spectral probes of the surprisal ledger}
% ============================================================

Let \(\eta\) be a probability measure on \(\R\) with finite second
moment
\[
\sigma_\eta^2
=
\int t^2\,\dd\eta(t)
<\infty.
\tag{36.1}
\]

Define
\[
\boxed{
D_\eta(P,Q)^2
=
\int
|\Phi_P(t)-\Phi_Q(t)|^2
\,\dd\eta(t).
}
\tag{36.2}
\]

For a translation-invariant positive-definite kernel whose Bochner
spectral measure is \(\eta\), this is the standard Fourier form of an
MMD-type discrepancy.
The novelty claimed here is not the metric construction itself, but its
use as an orbit-selected certificate inside the present counterfactual
source-cost bound.

\begin{theorem}[Spectral probe inequality]
\[
\boxed{
D_\eta(P,Q)
\le
\sigma_\eta
W_1(\mu_P,\mu_Q).
}
\tag{36.3}
\]
Hence
\[
\boxed{
W_1(\mu_P,\mu_Q)
\ge
\frac{D_\eta(P,Q)}{\sigma_\eta}.
}
\tag{36.4}
\]
\end{theorem}

\begin{proof}
For any coupling \((U,V)\) of \(\mu_P,\mu_Q\),
\[
|\Phi_P(t)-\Phi_Q(t)|
\le
\E|e^{itU}-e^{itV}|
\le
|t|\E|U-V|.
\]
Take the optimal \(W_1\) coupling, square, and integrate.
\end{proof}


% ============================================================
\section{Orbit--ledger--counterfactual bridge}
% ============================================================

Combining the spectral probe inequality with the synchronous
information-spectrum converse yields the main bridge.

\begin{theorem}[Orbit--ledger--counterfactual bridge]
Let \(P,Q\) be finite action distributions and let \(\eta\) be any
probability measure with
\[
0<\sigma_\eta^2<\infty.
\]
Then
\[
\boxed{
\CinfSyn(P,Q)
\ge
\frac{H(P)+H(Q)}2
+
\frac{D_\eta(P,Q)}{2\sigma_\eta}.
}
\tag{37.1}
\]

If \(\eta=\nu_{\psi,A}\) is the spectral measure of a normalized vector
\(\psi\) under a one-parameter unitary orbit generated by \(A\), then the
same inequality turns that geometric orbit into a valid spectral auditor
of the surprisal mismatch.
\end{theorem}

Thus
\[
\boxed{
\begin{gathered}
\text{one geometric vector }\psi
\\
\Downarrow
\\
\text{one-parameter orbit }e^{itA}\psi
\\
\Downarrow
\\
\text{orbit spectral measure }\eta
\\
\Downarrow
\\
D_\eta(P,Q)
\\
\Downarrow
\\
\text{synchronous source-cost certificate}.
\end{gathered}
}
\tag{37.2}
\]

No identification of the geometric kernel with a counterfactual latent
seed is required.
The bridge operates through common spectral geometry.


% ============================================================
\section{REAL kernels as spectral auditors}
% ============================================================

For \(k\ge1\), let \(\eta=\nu_k\) be the Mellin/dilation spectral measure
of the normalized polynomial kernel.
Then
\[
\sigma_{\nu_k}^2
=
\mathcal E_k
=
\frac{2k+1}{4(2k-1)}.
\]

Therefore
\[
\boxed{
\CinfSyn(P,Q)
\ge
\frac{H(P)+H(Q)}2
+
\frac{
D_{\nu_k}(P,Q)
}{
2\sqrt{\mathcal E_k}
}.
}
\tag{38.1}
\]

The parameter \(k\) selects a spectral auditor.

For \(k=0\), the second spectral moment diverges, so
(38.1) is unavailable.
Exactly in this singular case, however, the stronger direct identity
\[
2F_0(1/p,1/q)=\min(p,q)
\tag{38.2}
\]
recovers the hard \(W_1\) converse.

The family therefore exhibits a structural transition:
\[
\boxed{
\begin{array}{ccl}
\text{hard edge}
&\Longrightarrow&
\text{heavy Mellin tail}
\Longrightarrow
\text{linear cusp},
\\[1mm]
\text{finite dilation energy}
&\Longrightarrow&
\text{finite second spectral moment}
\Longrightarrow
\text{quadratic local resolution}.
\end{array}
}
\tag{38.3}
\]


% ============================================================
\section{Fractional-edge universality: the abstract target}
% ============================================================

The preceding examples motivate a general local classification.

Suppose an orbit correlation satisfies
\[
1-\Re R_{\psi,A}(t)
\sim
c|t|^\beta
\qquad
(t\to0)
\tag{39.1}
\]
for some
\[
0<\beta\le2.
\]

Define the associated normalized reciprocal-scale cost
\[
c_{\psi,A}(u,v)
=
\frac{u+v}{2}
-
\log_2
\Re R_{\psi,A}
\bigl(
(\ln2)|u-v|
\bigr),
\tag{39.2}
\]
whenever the correlation is positive near the origin.

Then the elementary expansion
\[
-\log(1-z)\sim z
\]
suggests
\[
\boxed{
c_{\psi,A}(u,v)
-
\frac{u+v}{2}
\sim
c(\ln2)^{\beta-1}
|u-v|^\beta.
}
\tag{39.3}
\]

The hard box corresponds to \(\beta=1\);
finite-energy smooth kernels correspond to \(\beta=2\).

A complete theorem should state the weakest positivity and regularity
conditions under which such an orbit profile induces a globally valid
probability-overlap functional and an operational counterfactual
converse.
The local asymptotic implication itself is elementary; the nontrivial
part is the global interface.


% ============================================================
\section{What the unified algebraic picture actually says}
% ============================================================

The phrase ``one vector plus an algebra'' has a precise but limited
meaning here.

For a finite classical source,
\[
\psi_\Gamma
\]
encodes the probability weights and
\[
\cA_\Gamma
\]
contains all classical observables.

For one observable generator \(A\), the vector state selects a spectral
measure
\[
\nu_{\psi,A}(B)
=
\ip{\psi}{E_A(B)\psi}.
\]

The complete one-parameter ledger is
\[
\ip{\psi}{e^{itA}\psi}
=
\int e^{it\xi}\,\dd\nu_{\psi,A}(\xi).
\]

The three principal uses in this paper are

\[
\boxed{
\begin{array}{c|c|c}
\text{setting}
&
\text{generator}
&
\text{vector spectral measure}
\\
\hline
\text{classical latent source}
&
M_f
&
\text{law of }f
\\
\text{surprisal ledger}
&
L_P
&
\mu_P
\\
\text{scale orbit}
&
A_{\mathrm{dil}}
&
\nu_K
\end{array}
}
\tag{40.1}
\]

This does not imply that the physical world is literally one vector.
Every finite classical probability space already admits such a
representation.

The substantive result is elsewhere:
the realization contract determines which observables must be
identified, and that algebraic identification changes the minimum
Shannon entropy of the latent source.


% ============================================================
\section{The meaning of the bill after Version 5}
% ============================================================

For a finite two-action memoryless channel, the causal theorem says
\[
\CinfCausal
=
\max_aH(P_a).
\tag{41.1}
\]

Thus unrestricted history-indexed causal coherence can asymptotically
push all additional counterfactual structure into a sublinear residual.

Stable synchronous identity can forbid that hiding mechanism.
For equal-entropy actions,
\[
\CinfSyn
\ge
h+\frac12W_1(\mu_P,\mu_Q).
\tag{41.2}
\]

The distinction can now be read in three equivalent languages:

\[
\boxed{
\begin{array}{ccl}
\text{probabilistic}
&:&
\text{same output law versus same latent event},
\\
\text{algebraic}
&:&
\text{history-dependent versus identified projections},
\\
\text{spectral}
&:&
\text{mean surprisal versus full ledger mismatch}.
\end{array}
}
\tag{41.3}
\]

For the crash test,
\[
\boxed{
\text{causal contract pays }1.5,
}
\]
while
\[
\boxed{
\text{synchronous stable identity pays at least }
1.728465468745531.
}
\]

The gap is therefore not a universal ``counterfactual tax''.
It is a contract-specific price of preserving a stronger identity
relation among unrealized events.


% ============================================================
\section{What is proved in Version 5}
% ============================================================

The following statements are established within the stated finite
models.

\begin{enumerate}[leftmargin=8mm]

\item C3-S policy safety is equivalent to correctness of every complete
action-string transversal.

\item The synchronous block cost is subadditive and admits a well-defined
asymptotic rate.

\item The synchronous transversal constraint rank is
\[
d_{\mathrm{obs}}^n.
\]

\item Hidden cross-time counterfactual synergy exists and an explicit
23-atom crash-test witness gives
\[
\Csyn_2
\le
3.878295576158564.
\]

\item The two-action synchronous information-spectrum converse is
\[
\Csyn_n
\ge
\frac n2
[
H(P)+H(Q)+W_1(\mu_P,\mu_Q)
].
\]

\item The converse extends to arbitrary finite action families through
the quantile envelope.

\item C3-C admits the action-tree quotient and the exact Bellman--MEC
recursion.

\item For equal-entropy actions,
\[
\Ccausal_n=nh+E_n
\]
with \(E_n\) nonnegative, nondecreasing, and subadditive.

\item For the crash test,
\[
E_2>E_1>0.
\]

\item For every finite two-action channel,
\[
\boxed{
\CinfCausal
=
\max\{H(P),H(Q)\}.
}
\]

\item The anti-diagonal construction yields
\[
\Ccausal_n
\le
nh_{\max}+O(n^{2/3}),
\]
and equal varentropy improves this to \(O(\sqrt n)\) residual.

\item For the crash test,
\[
\boxed{
\CinfCausal=1.5<\CinfSyn.
}
\]

\item Every finite classical latent source admits the one-vector
commutative-algebra representation
\[
(\psi_\Gamma,\cA_\Gamma).
\]

\item Stable synchronous event identity is exactly an identification of
event projections across histories.

\item The surprisal law is the vector spectral measure of the diagonal
surprisal operator.

\item
\[
d_{\mathrm{Lip}}(P,Q)=W_1(\mu_P,\mu_Q),
\]
so the synchronous converse is a lower bound in the distance between
two spectral states.

\item Arbitrarily many matching surprisal moments do not remove a
possible positive synchronous--causal gap; the dyadic parity ladder
provides an explicit family.

\item The polynomial scale kernels admit the explicit Mellin spectral
density (31.2), with a hard-edge Cauchy case at \(k=0\) and
finite-energy smooth cases for \(k\ge1\).

\item The reciprocal-scale \(k=0\) identity reproduces
\[
\min(p,q)
\quad\leftrightarrow\quad
\max(u,v),
\]
while higher \(k\) produce softer overlap costs.

\item Every finite-second-moment orbit spectral measure produces a valid
spectral-probe lower bound on the synchronous source cost through
Theorem~37.1.

\end{enumerate}


% ============================================================
\section{What is not proved}
% ============================================================

The following claims are not established.

\begin{enumerate}[leftmargin=8mm]

\item The 23-atom synchronous witness is not proved globally optimal.

\item The exact value of the crash-test synchronous asymptotic rate is
unknown.

\item The \(O(n^{2/3})\) or \(O(\sqrt n)\) finite-horizon causal residual
bounds are not claimed optimal.

\item The causal Shannon-floor theorem is proved here only for two
actions.

\item The finite-horizon anti-diagonal construction does not by itself
construct one projectively consistent infinite causal substrate that is
simultaneously optimal at every horizon.

\item No claim is made that the one-vector representation supplies a
physical ontology.
It is a classical representation theorem.

\item The weighted spectral discrepancy \(D_\eta\) is an MMD-type
construction, not a newly invented probability metric.

\item Mellin diagonalization of dilation and the spectral theorem are
standard.

\item Higher-\(k\) polynomial kernels do not strengthen the hard
\(k=0\) synchronous converse by the direct overlap argument.

\item No general theorem yet converts arbitrary fractional orbit-edge
exponents into sharp counterfactual source-cost exponents.

\item No historical priority claim is made for the algebraic,
Mellin-transform, MMD, Zolotarev, or operator-theoretic ingredients.

\end{enumerate}


% ============================================================
\section{Open problems after Version 5}
% ============================================================

\begin{problem}[Exact synchronous rate]
Determine the exact value of
\[
\CinfSyn
\]
for the crash-test pair.
\end{problem}

\begin{problem}[Exact finite causal residual]
Determine the true order of
\[
E_n=\Ccausal_n-nh.
\]
Is it bounded, logarithmic, a fractional power, or of another sublinear
order?
\end{problem}

\begin{problem}[Multi-action causal Shannon floor]
For
\[
|\cA|>2,
\]
does
\[
\CinfCausal(P_a:a\in\cA)
=
\max_aH(P_a)
\]
still hold?
\end{problem}

\begin{problem}[Projectively consistent infinite substrate]
Can the finite-horizon causal Shannon-floor construction be replaced by
one infinite projectively consistent causal response tree with an
appropriate entropy rate equal to \(h_{\max}\)?
\end{problem}

\begin{problem}[Sharp orbit auditor]
Among finite-energy orbit measures \(\eta\), determine those that maximize
\[
\frac{D_\eta(P,Q)}{\sigma_\eta}
\]
for a prescribed pair \(P,Q\).
How close can this spectral family come to the full \(W_1\) certificate?
\end{problem}

\begin{problem}[Fractional edge theorem]
Classify orbit vectors for which
\[
1-\Re R_{\psi,A}(t)
\asymp|t|^\beta
\]
and determine which such profiles induce globally valid overlap
functionals and counterfactual converses.
\end{problem}

\begin{problem}[Arithmetic versus analytic smoothing]
Find a resolution metric that separates low-frequency cumulant mismatch
from lattice-phase mismatch and admits both an ideal convolution law and
a minimum-entropy coupling interface.
\end{problem}

\begin{problem}[Higher-order counterfactual bridge]
Can Zolotarev-type or Fourier-weighted higher-order distances produce
stronger finite-horizon causal or synchronous bounds when lower moments
match?
\end{problem}

\begin{problem}[Algebraic contract classification]
Characterize source cost directly from relations imposed among
history-indexed event projections.
In particular, interpolate between complete C3-S identification and
fully history-indexed C3-C by partial equivalence relations on event
operators.
\end{problem}


% ============================================================
\section{Conclusion}
% ============================================================

The original accounting question began with a one-step observation:
if several mutually exclusive actions must already possess jointly
defined outputs, the minimum latent entropy can exceed the entropy of
any one observable action.

That statement survived, but its naive repeated interpretation did not.

Synchronous blocks can hide large amounts of dependency invisible to
every selectable transversal.
Nevertheless, the synchronous information-spectrum converse shows that
stable cross-history event identity can retain a strictly positive
asymptotic cost when the self-information spectra differ.

General causal trees behave differently.
By placing correlations beyond the causal anti-diagonal, a two-action
history-indexed realization asymptotically reaches the ordinary Shannon
floor.

The resulting crash-test hierarchy is

\[
\boxed{
1.5
=
\CinfCausal
<
1.728465468745531
\le
\CinfSyn
\le
1.939147788079282
<
2.053918585607260
=
C_1.
}
\tag{46.1}
\]

Version 5 shows that this separation also has a compact algebraic form.

A classical latent law is one vector
\[
\psi
\]
together with a commutative algebra of observables.

The difference between C3-C and C3-S is an algebraic difference in which
event projections are identified across histories.

The self-information of an action is one diagonal operator
\[
L_P.
\]
Its first spectral moment is Shannon entropy.
Its second centered moment is varentropy.
Its full vector spectral measure is the surprisal law.
Its full one-parameter characteristic function is the spectral ledger.

The causal two-action rate depends on the largest first moment.
The synchronous converse can depend on the distance between the complete
spectral ledgers.

The scale-kernel calculation fits the same template.
One geometric vector under a unitary dilation orbit generates one
Mellin spectral measure.
That measure can be used as a spectral auditor of the surprisal ledger
and therefore as a lower-bound certificate on synchronous source cost.

What remains from the earliest intuition is therefore not an ontological
claim about a privileged physical vector.

It is a mathematical architecture:

\begin{center}
\fbox{
\begin{minipage}{0.84\textwidth}
\centering
\textbf{One state vector, an algebra acting on it, and a spectral ledger
recording what that algebra can resolve.}
\end{minipage}}
\end{center}

The bill is determined not only by what can happen, but by which
unrealized events the realization contract insists must remain
identically the same.


% ============================================================
\appendix
% ============================================================

\section{Explicit 23-atom synchronous witness}
\label{app:witness}

Write each atom as
\[
(x_1,z_1,x_2,z_2)\in\{0,1,2\}^4.
\]

The nonzero probabilities are

\[
\renewcommand{\arraystretch}{1.32}
\begin{array}{c|c}
(x_1,z_1,x_2,z_2) & \widetilde G(x_1,z_1,x_2,z_2)
\\
\hline
(0,0,0,0) & \dfrac{(2q-1)(12q-5)}8 \\
(0,0,0,2) & -q(2q-1) \\
(0,0,2,1) & \dfrac q4 \\
(0,1,0,0) & q^2 \\
(0,1,0,2) & -\dfrac{(2q-1)(8q-3)}8 \\
(0,1,1,1) & \dfrac q4 \\
(0,2,1,0) & -\dfrac{2q-1}{8} \\
(0,2,2,0) & -\dfrac{(2q-1)(8q-3)}4 \\
(0,2,2,2) & \dfrac{(2q-1)(16q-7)}8 \\
(1,0,0,0) & \dfrac q4 \\
(1,0,1,1) & \dfrac{4q-1}{16} \\
(1,1,0,2) & -\dfrac{2q-1}{8} \\
(1,1,2,1) & \dfrac1{16} \\
(1,2,1,2) & -\dfrac{2q-1}{8} \\
(2,0,0,0) & -\dfrac{32q^2-40q+11}{16} \\
(2,0,0,1) & \dfrac{(4q-1)^2}{16} \\
(2,0,1,0) & \dfrac1{16} \\
(2,1,0,1) & \dfrac{16q^2-4q-1}{16} \\
(2,1,0,2) & -\dfrac{12q-5}{16} \\
(2,1,2,2) & \dfrac{4q-1}{16} \\
(2,2,0,0) & \dfrac{(2q-1)^2}{2} \\
(2,2,0,1) & -q(2q-1) \\
(2,2,2,0) & -\dfrac{2q-1}{8}
\end{array}
\tag{A.1}
\]

At the root \(q\) from (5.2), these values are positive, sum to
one, and each of the four selected coordinate pairs has exactly the
required product law.


% ============================================================
\section{One-shot spectral rigidity}
\label{app:rigidity}

The five one-shot vertex-spectrum types may be represented by

\[
\lambda_\star
=
\left(
q,\frac14,1-2q,\frac12-q,2q-\frac34
\right),
\]

\[
\lambda_A
=
\left(
\frac14,\frac14,1-2q,q-\frac14,q-\frac14
\right),
\]

\[
\lambda_B
=
\left(
\frac34-q,\frac14,1-2q,q-\frac14,2q-\frac34
\right),
\]

\[
\lambda_C
=
\left(
2q-\frac12,\frac14,1-2q,q-\frac14,\frac12-q
\right),
\]

and

\[
\lambda_D
=
\left(
q,\frac14,q-\frac14,\frac12-q,\frac12-q
\right).
\]

On
\[
\frac38<q<\frac5{12},
\]
the cumulative-sum differences between \(\lambda_\star\) and the other
types are

\[
\begin{array}{c|ccccc}
&
1&2&3&4&5
\\
\hline
A&
q-\frac14&
q-\frac14&
q-\frac14&
\frac12-q&
0
\\[1mm]
B&
2q-\frac34&
2q-\frac34&
2q-\frac34&
0&
0
\\[1mm]
C&
\frac12-q&
\frac12-q&
\frac12-q&
\frac54-3q&
0
\\[1mm]
D&
0&
0&
\frac54-3q&
\frac54-3q&
0 .
\end{array}
\tag{B.1}
\]

All displayed nonzero entries are positive.
Thus \(\lambda_\star\) strictly majorizes every competing type, and
strict Schur concavity of Shannon entropy proves one-shot spectrum
rigidity.


% ============================================================
\section{Reproducibility ledger}
% ============================================================

The numerical values used in the crash test are

\[
q
=
0.41003242818198557608\ldots,
\]
\[
C_1
=
2.0539185856072597\ldots,
\]
\[
H(\widetilde G)
=
3.878295576158564\ldots,
\]
\[
\frac12H(\widetilde G)
=
1.939147788079282\ldots,
\]
\[
W_1(\mu_P,\mu_Q)
=
0.4569309374910615\ldots,
\]
and
\[
1.5+\frac12W_1
=
1.7284654687455308\ldots.
\]

The Version 4 companion script checks the equal-entropy root, the
one-shot witness, the 23-atom synchronous law, all four transversal
identities, the information-spectrum lower bound, one finite
anti-diagonal causal realization by complete state enumeration, and the
entropy formula for that source.

Version 5 adds analytic operator and Mellin identities.
A dedicated symbolic/numerical audit of the explicit spectral formulas
(31.1)--(34.9) is recommended as the next companion
reproducibility layer.


% ============================================================
\section{Literature status table}
% ============================================================

\begin{center}
\renewcommand{\arraystretch}{1.25}
\begin{tabular}{
>{\raggedright\arraybackslash}p{0.40\textwidth}
>{\raggedright\arraybackslash}p{0.18\textwidth}
>{\raggedright\arraybackslash}p{0.32\textwidth}
}
\toprule
Topic & Status used here & Representative reference
\\
\midrule

Minimum-entropy coupling
&
Known
&
Kovačević--Stanojević--Šenk
\\

Information-spectrum MEC converse
&
Known
&
Shkel--Yadav
\\

Approximate multi-marginal MEC
&
Known
&
Li
\\

Profile MEC bound
&
Known
&
Compton et al.\ 2023
\\

PTAS for constant-number MEC
&
Known, 2026
&
Compton
\\

Exact common information / synthesis
&
Known
&
Kumar--Li--El Gamal; Yu--Tan
\\

Online randomness recycling
&
Known, active
&
Draper--Saad et al.
\\

Event-keyed common randomness
&
Known neighboring issue
&
Buffalo--Pearson--Klein
\\

Mellin analysis of scale invariance
&
Known
&
Gupta--Singh--Aggarwal--Joshi
\\

MMD and spectral distribution embeddings
&
Known
&
Sriperumbudur et al.; Gretton et al.
\\

Group-invariant MMD
&
Known recent direction
&
Giacofci--Meynaoui--Podgorny
\\

Fourier/Zolotarev higher-order metrics
&
Known
&
Bobkov--Ledoux and classical literature
\\

Synchronous all-transversal source rate
&
Not located in this formulation
&
Compare MEC and structured-independence literature
\\

Two-action anti-diagonal causal Shannon floor
&
Derived here
&
No direct precedent located in present search
\\

Explicit REAL Mellin spectrum and counterfactual auditor bridge
&
Derived here
&
Uses standard Mellin, spectral-theorem, and MMD ingredients
\\

C3-S/C3-C algebraic operator-identification formulation
&
Derived here
&
Compare classical commutative operator representation
\\

\bottomrule
\end{tabular}
\end{center}


% ============================================================
\begin{thebibliography}{99}
% ============================================================

\bibitem{Kovacevic2015}
M. Kovačević, I. Stanojević, and V. Šenk,
\emph{On the Entropy of Couplings},
Information and Computation 242 (2015), 369--382.
\url{https://arxiv.org/abs/1303.3235}

\bibitem{ShkelYadav2023}
Y. Y. Shkel and A. K. Yadav,
\emph{Information Spectrum Converse for Minimum Entropy Couplings and
Functional Representations},
2023 IEEE International Symposium on Information Theory (ISIT), 2023.
\url{https://arxiv.org/abs/2305.05745}

\bibitem{Li2021}
C. T. Li,
\emph{Efficient Approximate Minimum Entropy Coupling of Multiple
Probability Distributions},
IEEE Transactions on Information Theory 67(8) (2021), 5259--5279.
\url{https://arxiv.org/abs/2006.07955}

\bibitem{Compton2023}
S. Compton, D. Katz, B. Qi, K. Greenewald, and M. Kocaoglu,
\emph{Minimum-Entropy Coupling Approximation Guarantees Beyond the
Majorization Barrier},
Proceedings of AISTATS 2023,
PMLR 206 (2023), 10445--10469.
\url{https://proceedings.mlr.press/v206/compton23a.html}

\bibitem{Compton2026PTAS}
S. Compton,
\emph{Efficient $\varepsilon$-Approximate Minimum-Entropy Couplings},
IEEE Transactions on Information Theory, 2026.
doi:10.1109/TIT.2026.3710290.
Preprint:
\url{https://arxiv.org/abs/2509.19598}

\bibitem{KumarLiElGamal}
G. R. Kumar, C. T. Li, and A. El Gamal,
\emph{Exact Common Information},
2014.
\url{https://arxiv.org/abs/1402.0062}

\bibitem{YuTanExact}
L. Yu and V. Y. F. Tan,
\emph{Exact Channel Synthesis},
2018.
\url{https://arxiv.org/abs/1810.13246}

\bibitem{DraperSaad2026}
T. L. Draper and F. A. Saad,
\emph{Efficient Online Random Sampling via Randomness Recycling},
Proceedings of the 2026 Annual ACM--SIAM Symposium on Discrete
Algorithms (SODA), 2473--2511, 2026.
doi:10.1137/1.9781611978971.89.

\bibitem{DraperHarrisSaad2026}
T. L. Draper, D. G. Harris, and F. A. Saad,
\emph{Online Random Sampling with Real Probabilities},
arXiv:2607.13828, 2026.

\bibitem{DraperSaadSpace2026}
T. L. Draper and F. A. Saad,
\emph{Space--Entropy Lower Bounds for Random Sampling},
arXiv:2607.14503, 2026.

\bibitem{BuffaloPearsonKlein2026}
V. Buffalo, C. A. B. Pearson, and D. Klein,
\emph{Realizing Common Random Numbers:
Event-Keyed Hashing for Causally Valid Stochastic Models},
arXiv:2603.11084, 2026.

\bibitem{GuptaSinghAggarwalJoshi2025}
A. Gupta, P. Singh, P. Aggarwal, and S. D. Joshi,
\emph{Unified Framework for Linear Scale Invariant Signals, Systems,
and Transforms: A Tutorial},
Digital Signal Processing 157 (2025), 104880.
doi:10.1016/j.dsp.2024.104880.

\bibitem{Sriperumbudur2010}
B. K. Sriperumbudur, A. Gretton, K. Fukumizu, B. Schölkopf,
and G. R. G. Lanckriet,
\emph{Hilbert Space Embeddings and Metrics on Probability Measures},
Journal of Machine Learning Research 11 (2010), 1517--1561.

\bibitem{Gretton2012}
A. Gretton, K. M. Borgwardt, M. J. Rasch, B. Schölkopf, and A. Smola,
\emph{A Kernel Two-Sample Test},
Journal of Machine Learning Research 13 (2012), 723--773.

\bibitem{GiacofciMeynaouiPodgorny2026}
M. Giacofci, A. Meynaoui, and A. Podgorny,
\emph{A Kernel Two-Sample Test Invariant under Group Action with
Applications to Functional Data},
arXiv:2603.16294, 2026.

\bibitem{BobkovLedoux2026}
S. G. Bobkov and M. Ledoux,
\emph{Fourier Analytic Bounds for Zolotarev Distances and Applications
to Empirical Measures},
in \emph{Progress in Probability}, Vol. 82,
Birkhäuser, 2026, pp. 259--312.
\url{https://doi.org/10.1007/978-3-032-06057-0_9}.

\bibitem{Kuhn1953}
H. W. Kuhn,
\emph{Extensive Games and the Problem of Information},
in H. W. Kuhn and A. W. Tucker (eds.),
\emph{Contributions to the Theory of Games II},
Annals of Mathematics Studies 28,
Princeton University Press, 1953.

\bibitem{MainRandour2024}
J. C. A. Main and M. Randour,
\emph{Different Strokes in Randomised Strategies:
Revisiting Kuhn's Theorem under Finite-Memory Assumptions},
Information and Computation 301 (2024), 105229.
\url{https://arxiv.org/abs/2201.10825}

\bibitem{ContLim}
R. Cont and F. R. Lim,
\emph{Causal Transport on Path Space},
arXiv:2412.02948, 2024.

\bibitem{GavinskyPudlak}
D. Gavinsky and P. Pudlák,
\emph{On the Joint Entropy of $d$-Wise-Independent Variables},
arXiv:1503.08154, 2015.

\bibitem{KahleWenzelAy}
T. Kahle, W. Wenzel, and N. Ay,
\emph{Hierarchical Models, Marginal Polytopes, and Linear Codes},
Kybernetika 45(2) (2009), 189--207.
\url{https://arxiv.org/abs/0805.1301}

\bibitem{RassouliRosasGunduz}
B. Rassouli, F. E. Rosas, and D. Gündüz,
\emph{Data Disclosure under Perfect Sample Privacy},
IEEE Transactions on Information Forensics and Security 15 (2020),
2012--2025.
\url{https://arxiv.org/abs/1904.01711}

\bibitem{Wilming2021}
H. Wilming,
\emph{Entropy and Reversible Catalysis},
Physical Review Letters 127 (2021), 260402.
\url{https://arxiv.org/abs/2012.05573}

\bibitem{Wilming2022}
H. Wilming,
\emph{Correlations in Typicality and an Affirmative Solution to the
Exact Catalytic Entropy Conjecture},
Quantum 6 (2022), 858.
\url{https://arxiv.org/abs/2205.08915}

\bibitem{KeaneSmorodinsky}
M. Keane and M. Smorodinsky,
\emph{Bernoulli Schemes of the Same Entropy Are Finitarily Isomorphic},
Annals of Mathematics 109(2) (1979), 397--406.

\bibitem{HarveyPeres}
N. Harvey and Y. Peres,
\emph{An Invariant of Finitary Codes with Finite Expected Square Root
Coding Length},
Ergodic Theory and Dynamical Systems 31(1) (2011), 77--90.
\url{https://arxiv.org/abs/math/0309120}

\bibitem{Gabor2024}
U. Gabor,
\emph{Moment-Optimal Finitary Isomorphism for i.i.d. Processes of Equal
Entropy},
arXiv:2412.15425, 2024.

\end{thebibliography}

\end{document}
